\subsection{Problems with RNNs: Vanishing and Exploding Gradients}
\label{sec:vanishing-exploding-gradients}
\slidesrc{Language models \& RNNs: 85--101}

In practice, training the RNN presented so far (often called a \textbf{vanilla RNN},
i.e.\ the plain basic version) is hard, because of two related problems:
\begin{enumerate}[nosep]
  \item \textbf{vanishing gradients};
  \item \textbf{exploding gradients}.
\end{enumerate}
Both come from the same source: in BPTT the gradient travels backwards through many
steps, and at each step it is \emph{multiplied} by the same matrix.

\needspace{12\baselineskip}
\subsubsection{Vanishing gradients: intuition}
\slidesrc{Language models \& RNNs: 87--91}

How much does the loss at step 4 depend on the hidden state at step 1? By the chain
rule, the influence of $\vect h^{(1)}$ on $J^{(4)}$ passes through $\vect h^{(2)}$,
$\vect h^{(3)}$ and $\vect h^{(4)}$, and each passage contributes one factor:
\[
  \frac{\partial J^{(4)}}{\partial \vect h^{(1)}} =
  \frac{\partial \vect h^{(2)}}{\partial \vect h^{(1)}} \times
  \frac{\partial \vect h^{(3)}}{\partial \vect h^{(2)}} \times
  \frac{\partial \vect h^{(4)}}{\partial \vect h^{(3)}} \times
  \frac{\partial J^{(4)}}{\partial \vect h^{(4)}} .
\]

\begin{center}
\begin{tikzpicture}[>=Stealth,
    neu/.style={circle, draw=cyan!70!black, fill=cyan!35, minimum size=5pt, inner sep=0pt},
    hid/.style={draw=cyan!70!black, rounded corners, thick, minimum width=0.5cm,
                minimum height=1.4cm}]
  \foreach \i in {1,...,4} {
    \node[hid] (h\i) at (2.6*\i,0) {};
    \foreach \k in {1,...,4} \node[neu] at (2.6*\i,0.7-0.28*\k) {};
    \node[font=\small, anchor=north] at (h\i.south) {$\vect h^{(\i)}$};
  }
  \foreach \i/\j in {1/2,2/3,3/4} \draw[->, nbgray] ([yshift=4pt]h\i.east) --
      node[above, font=\scriptsize] {$W_h$} ([yshift=4pt]h\j.west);
  \node[font=\small, nbred] (J) at (10.4,1.9) {$J^{(4)}$};
  \draw[->, nbgray] ([xshift=-3pt]h4.north) -- ([xshift=-3pt]J.south);
  % backward arrows getting thinner
  \draw[->, nbred, line width=3pt] ([xshift=4pt]J.south) -- ([xshift=4pt]h4.north);
  \draw[->, nbred, line width=2.2pt] ([yshift=-6pt]h4.west) -- ([yshift=-6pt]h3.east);
  \draw[->, nbred!80, line width=1.3pt] ([yshift=-6pt]h3.west) -- ([yshift=-6pt]h2.east);
  \draw[->, nbred!60, line width=0.5pt] ([yshift=-6pt]h2.west) -- ([yshift=-6pt]h1.east);
  \node[font=\scriptsize, nbred, align=center] at (5.2,-1.45)
       {the gradient shrinks at every step backwards};
\end{tikzpicture}
\captionof{figure}{The gradient of $J^{(4)}$ flows backwards from $\vect h^{(4)}$ to
$\vect h^{(1)}$ (red arrows). If each factor is ``small'', the signal reaching the first
tokens is tiny.}
\end{center}

\begin{nbintuition}[multiplying small numbers]
\notinslides
If each factor were $0.5$, after 3 steps the signal would be $0.5^3 = 0.125$, after 10
steps $0.001$, after 20 steps one millionth. A product of many numbers smaller than 1
\emph{vanishes}; a product of many numbers larger than 1 \emph{explodes}. It is the
same as the children's game ``telephone'': a message whispered along a long chain of
people arrives at the end faint or distorted.
\end{nbintuition}

\begin{nbdef}{Vanishing gradients}
When the factors $\partial \vect h^{(t)} / \partial \vect h^{(t-1)}$ are small, the
gradients with respect to the \textbf{first tokens become tiny}: the loss at a late step
hardly influences how the network processes early words.
\end{nbdef}

\subsubsection{Vanishing gradients: proof sketch}
\slidesrc{Language models \& RNNs: 92--95}

Use again the simplified model $\vect h^{(t)} = W_h \vect h^{(t-1)}$,
$J^{(t)}(\theta) = \mathrm{CE}(\vect h^{(t)}, \vect y^{*(t)})$. Consider the loss at step
$i$ and its gradient with respect to the hidden state of an earlier step $j < i$. Call
$l = i - j$ the \textbf{distance} between the two steps.
\begin{align*}
  \frac{\partial J^{(i)}(\theta)}{\partial \vect h^{(j)}}
  &= \frac{\partial J^{(i)}(\theta)}{\partial \vect h^{(i)}}
     \prod_{j < t \le i} \frac{\partial \vect h^{(t)}}{\partial \vect h^{(t-1)}}
     \reason{chain rule, one factor per step}\\
  &= \frac{\partial J^{(i)}(\theta)}{\partial \vect h^{(i)}}
     \prod_{j < t \le i} W_h
     \reason{$\vect h^{(t)} = W_h \vect h^{(t-1)}$}\\
  &= \frac{\partial J^{(i)}(\theta)}{\partial \vect h^{(i)}}\, W_h^{\,l}
     \reason{$l$ equal factors}
\end{align*}
If $W_h$ is ``small'', the power $W_h^{\,l}$ becomes minuscule as $l$ grows. What
``small'' means for a matrix is made precise by its \textbf{eigenvalues}.

\notinslides
Recall that an eigenvector $\vect q$ of a matrix is a direction that the matrix only
stretches: multiplying by the matrix gives $\lambda \vect q$, where the number
$\lambda$ is the eigenvalue. Multiplying $l$ times stretches it by $\lambda^l$.

Suppose $W_h$ has eigenvalues $\lambda_1, \ldots, \lambda_n$ with
$|\lambda_k| < 1$ and (left) eigenvectors $\vect q_1, \ldots, \vect q_n$ that form a
basis. Write the gradient row vector in this basis,
$\frac{\partial J^{(i)}}{\partial \vect h^{(i)}} = \sum_k c_k \vect q_k\T$. Each
multiplication by $W_h$ multiplies the $k$-th component by $\lambda_k$, so
\[
  \frac{\partial J^{(i)}(\theta)}{\partial \vect h^{(i)}}\, W_h^{\,l}
  = \sum_{k=1}^{n} c_k\, \lambda_k^{\,l}\, \vect q_k\T
  \;\approx\; 0 \quad \text{for large } l ,
\]
because every $\lambda_k^{\,l} \to 0$: the gradient \textbf{vanishes}. How fast?

\begin{center}
\small
\begin{tabular}{@{}c|ccccc@{}}
\toprule
$\lambda$ & $l = 1$ & $l = 5$ & $l = 10$ & $l = 20$ & $l = 50$ \\
\midrule
0.5 & 0.5 & 0.031 & 0.00098 & $9.5 \times 10^{-7}$ & $\approx 0$ \\
0.9 & 0.9 & 0.59  & 0.35    & 0.12                 & 0.0052 \\
1.1 & 1.1 & 1.61  & 2.59    & 6.73                 & 117 \\
\bottomrule
\end{tabular}
\captionof{table}{$\lambda^l$: the factor by which a component of the gradient is
multiplied after travelling back $l$ steps. Below 1 it decays exponentially, above 1 it
grows exponentially.}
\end{center}

The same phenomenon occurs with more realistic RNN models (reference: Pascanu et al.,
\emph{On the difficulty of training recurrent neural networks}, 2013).

\begin{nbwarning}[details the slide leaves out]
\begin{itemize}[nosep]
  \item The condition is on the \emph{absolute value}: $|\lambda_k| < 1$. The slide
        writes $\lambda_k < 1$, but an eigenvalue $-2$ is smaller than 1 and still
        makes the gradient explode. (Eigenvalues can even be complex numbers; then
        $|\lambda|$ is their modulus.)
  \item The argument assumes $W_h$ has a basis of eigenvectors (it is
        \emph{diagonalisable}); the general case is similar, using the largest
        singular value instead.
  \item The slide uses $i$ both for the time step and for the eigenvalue index; above
        the eigenvalue index is $k$.
  \item In a real RNN, $\vect h^{(t)} = \sigma(W_h \vect h^{(t-1)} + \ldots)$, so each
        factor is $\mathrm{diag}\bigl(\sigma'(\cdot)\bigr) W_h$. The derivative of the
        sigmoid is at most $\frac14$ (of $\tanh$ at most 1), which makes the factors
        even smaller and vanishing even more likely.
\end{itemize}
\end{nbwarning}

\subsubsection{Why it matters: long-term dependencies}
\slidesrc{Language models \& RNNs: 96--98}

The gradient that updates the weights is the sum of contributions from all steps. A
contribution from a \textbf{nearby token} arrives almost intact; one from a
\textbf{far-away token} arrives multiplied by $W_h^{\,l}$, i.e.\ almost 0. So the
weights are updated almost only to exploit short-range patterns.

\begin{nbfact}{Long-term dependencies}
With vanishing gradients, \textbf{long-term dependencies are (practically) impossible
to learn}, because the gradient from far-away tokens is much smaller than from nearby
tokens.
\end{nbfact}

\begin{nbexample}{a dependency 20+ tokens back}
Predict the next word:
\begin{quote}\small
After leaving \textbf{work}, he went to pick up his daughter from school. They had
planned to go to the pet shop to buy a hamster and drove through town looking for one
she liked. Eventually they found it, but before heading back home he had to rush back
to his \underline{\quad\quad}
\end{quote}
The right answer is \emph{work}. To predict it the model must connect the blank with
the word \emph{work} at the very beginning, more than 20 tokens earlier. During
training, the gradient that would teach this connection is tiny compared with the
gradients of short-range patterns (``his'' $\to$ a noun, ``rush back to his'' $\to$
\emph{car}, \emph{house}, \ldots), so it provides little signal and learning
long-term dependencies becomes very slow. A vanilla RNN tends to guess from the last
few words only, i.e.\ it behaves much like an n-gram model.
\end{nbexample}

\subsubsection{Exploding gradients and gradient clipping}
\slidesrc{Language models \& RNNs: 99--100}

Symmetrically, if the eigenvalues of $W_h$ are larger than 1 (in absolute value), the
gradients become \textbf{too big}:
\begin{itemize}[nosep]
  \item an update $\theta \leftarrow \theta - \alpha \nabla_\theta J$ with a huge
        gradient is a huge step, which can jump over the minimum and
        \textbf{increase} the loss instead of decreasing it;
  \item in the worst case the numbers overflow and the weights become
        \texttt{inf} or \texttt{NaN} (``not a number''): training is ruined.
\end{itemize}

\begin{nbdef}{Gradient clipping}
If the gradient is too large, scale it down. With a threshold $\tau > 0$:
\begin{enumerate}[nosep]
  \item set $\hat{\vect g} \leftarrow \nabla_\theta J(\theta)$;
  \item \textbf{if} $\norm{\hat{\vect g}} \ge \tau$:
  \item \quad set $\hat{\vect g} \leftarrow \dfrac{\tau}{\norm{\hat{\vect g}}}\,
        \hat{\vect g}$;
  \item \textbf{end if}
\end{enumerate}
then update the weights with $\hat{\vect g}$.
\end{nbdef}

Multiplying by the positive number $\tau / \norm{\hat{\vect g}}$ does not change the
direction of the vector, only its length, which becomes exactly $\tau$. So
\textbf{the update direction remains the same, but the step size is reduced}.

\begin{nbexample}{clipping a gradient}
\notinslides
Gradient $\hat{\vect g} = (30, 40)$, threshold $\tau = 5$. Its length is
$\norm{\hat{\vect g}} = \sqrt{30^2 + 40^2} = 50 \ge 5$, so
\[
  \hat{\vect g} \leftarrow \frac{5}{50}\,(30, 40) = (3, 4),
\]
which has length $\sqrt{9 + 16} = 5 = \tau$ and points in the same direction (the
ratio between the entries is still $3 : 4$).
\end{nbexample}

\begin{center}
\begin{tikzpicture}[>=Stealth]
\begin{axis}[width=8.5cm, height=5.5cm, axis lines=left, xmin=-0.6, xmax=4.6, ymin=0, ymax=5,
    xlabel={parameter $\theta$}, ylabel={loss $J(\theta)$}, xtick=\empty, ytick=\empty,
    label style={font=\small}, clip=false]
  \addplot[domain=-0.3:4.3, samples=200, thick, nbblue]
     {0.15*(x-1)^2 + 3/(1+exp(-8*(x-3)))};
  \fill[black] (axis cs:3,2.1) circle (2pt);
  \draw[->, thick, nbgreen] (axis cs:3,2.1) -- (axis cs:2.5,0.39);
  \node[font=\scriptsize, nbgreen, align=left, anchor=west] at (axis cs:3.15,1.0)
       {with clipping:\\small step, same direction};
  \draw[->, thick, dashed, nbred] (axis cs:3,2.1) .. controls (axis cs:2.0,4.4) ..
       (axis cs:-0.2,4.1);
  \node[font=\scriptsize, nbred, align=center] at (axis cs:1.5,4.7)
       {without clipping: thrown far away};
  \node[font=\scriptsize, nbgray, anchor=west] at (axis cs:3.3,2.3) {cliff};
\end{axis}
\end{tikzpicture}
\captionof{figure}{A ``cliff'' in the loss, typical of RNNs. On the steep wall the
gradient is enormous: a plain gradient step throws the parameters far away, undoing
training; the clipped step goes in the same direction but stays local.}
\end{center}

The slide illustrates this with the 3D loss surface from Pascanu et al.: a flat region
with a steep wall, where unclipped updates bounce off the wall, while clipped ones
proceed smoothly. In practice \textbf{clipping essentially solves the exploding
gradient problem}. It does nothing, however, for vanishing gradients: one cannot
enlarge a gradient that carries no information.

\subsubsection{Solutions for vanishing gradients}
\slidesrc{Language models \& RNNs: 101}

Vanilla RNNs cannot learn long-range dependencies because \textbf{they do not have a
memory}: in
\[
  \vect h^{(t)} = \sigma\bigl(W_h \vect h^{(t-1)} + W_e \vect e^{(t)}\bigr)
\]
the hidden state is \textbf{constantly rewritten}: at every step all of it is
multiplied by $W_h$ and squashed again, so old information fades. The solutions
follow two ideas:
\begin{itemize}
  \item \textbf{Separate memory mechanisms}: \textbf{LSTMs} (Long Short-Term Memory) and
        \textbf{GRUs} (Gated Recurrent Units) add a memory cell and \emph{gates}, small
        learnt switches that decide at each step what to keep, what to forget and what
        to write. Information the gates decide to keep passes from step to step almost
        unchanged, so its gradient does not vanish. These were the
        \textbf{state-of-the-art NLP models in the 2010s}.
  \item \textbf{Direct connections between far-away tokens}: \textbf{attention} lets
        each position look directly at any earlier position, and \textbf{residual
        connections} add a shortcut that skips over layers. The gradient can then
        travel along a short path instead of through dozens of multiplications.
        Attention is the core of the transformer, the architecture of today's LLMs.
\end{itemize}

\begin{nbsummary}[vanishing and exploding gradients]
\begin{itemize}[nosep]
  \item Gradient across $l$ steps: $\frac{\partial J^{(i)}}{\partial \vect h^{(j)}} =
        \frac{\partial J^{(i)}}{\partial \vect h^{(i)}} W_h^{\,l}$ (simplified model).
  \item All $|\lambda_k| < 1$ $\Rightarrow$ $\lambda_k^l \to 0$: \textbf{vanishing}
        gradients, long-term dependencies are not learnt.
  \item Some $|\lambda_k| > 1$ $\Rightarrow$ \textbf{exploding} gradients: loss jumps
        up, \texttt{inf}/\texttt{NaN}.
  \item Exploding $\to$ \textbf{gradient clipping}: if $\norm{\vect g} \ge \tau$, use
        $\frac{\tau}{\norm{\vect g}}\vect g$ (same direction, length $\tau$).
  \item Vanishing $\to$ memory with gates (LSTM, GRU) or direct connections
        (attention, residual connections).
\end{itemize}
\end{nbsummary}

\begin{nbexercise}{3}{vanishing or exploding?}
\notinslides
In the simplified model $h^{(t)} = w\,h^{(t-1)}$ with a scalar hidden state, the
gradient reaching a token $l$ steps back is multiplied by $w^l$. For $w = 0.8$ and
$w = 1.25$, after how many steps does the factor fall below $\frac{1}{100}$ or rise
above $100$? What happens with $w = -1.25$?
\end{nbexercise}
\begin{nbsolution}
For $w = 0.8$ we need $0.8^l < 0.01$, i.e.\ $l > \frac{\ln 0.01}{\ln 0.8} =
\frac{-4.605}{-0.223} \approx 20.6$: from $l = 21$ on (check: $0.8^{21} \approx 0.0092$).

For $w = 1.25$ we need $1.25^l > 100$, i.e.\ $l > \frac{\ln 100}{\ln 1.25} =
\frac{4.605}{0.223} \approx 20.6$: again from $l = 21$ on ($1.25^{21} \approx 108$).
(The two answers coincide because $1.25 = 1/0.8$.)

For $w = -1.25$: $|w^l| = 1.25^l$, the same growth, with a sign that alternates at
every step. The gradient explodes even though $w < 1$, which is why the condition must
be stated on $|w|$ (on $|\lambda|$ for matrices).

Lesson: about 20 steps are enough to lose or blow up the signal by a factor 100, even
with a factor only 20\% away from 1. This is the same order of magnitude as the
truncation length used in BPTT and as the distance in the ``rush back to his work''
example.
\end{nbsolution}
